\section{Attack Implementations}
\label{sec:attacks}

\subsection{AgentPoison}

\agentpoison{} \citep{chen2024agentpoison} optimizes a trigger sequence
$\trigger^* = (t_1, \ldots, t_\ell)$ such that for any victim query
$q \in \mathcal{Q}_v$, the triggered embedding $E(q \oplus \trigger)$ is
close to the adversarial passage embedding $E(\passage)$:
\begin{equation}
  \trigger^* = \arg\max_\trigger\; \EE_{q \sim \mathcal{Q}_v}
  \left[\cos\!\left(E(q \oplus \trigger),\, E(\passage)\right)\right].
  \label{eq:agentpoison_obj}
\end{equation}

\paragraph{DPR HotFlip trigger optimization.}
We implement the paper-faithful gradient-guided trigger optimizer using the
DPR context encoder (\texttt{facebook/dpr-ctx\_encoder-single-nq-base},
768-dim BERT-based).  At each iteration, the gradient of the cosine-similarity
loss with respect to the trigger token embedding at position $i$ is computed:
\begin{equation}
  t_i^* = \arg\max_{v \in \mathcal{V}}\;
  \nabla_{e_{t_i}} \mathcal{L} \cdot e_v,
  \label{eq:hotflip}
\end{equation}
where $\mathcal{V}$ is the filtered BERT vocabulary (alphabetic, non-subword
tokens, $|\mathcal{V}| \approx 18{,}000$) and $e_v$ is the embedding of token
$v$.  Top-$n_{\text{cand}}=32$ candidates are evaluated with exact forward
passes; a GPT-2 perplexity gate ($\mathrm{ppl} < 500$) ensures trigger
fluency.  Figure~\ref{fig:dpr_convergence} shows convergence.

\paragraph{Centroid passage.}
The adversarial passage is built to target the centroid of victim query
embeddings:
\begin{equation}
  \passage = \mathrm{CentroidPassage}(\mathcal{Q}_v)
  = \text{passage covering vocabulary}\; \bigcup_{q \in \mathcal{Q}_v} \mathrm{KeyTerms}(q).
  \label{eq:centroid}
\end{equation}
$\mathrm{KeyTerms}$ extracts content words (stopword-filtered, frequency-ranked).
This places $\passage$ near the centroid of victim query embeddings, achieving
moderate-to-high cosine similarity with \emph{every} victim query.
The centroid passage plus DPR-optimized trigger achieves mean cosine
similarity 0.78 vs.\ 0.48 baseline (random trigger) and 0.71 (centroid only).

\paragraph{Triggered-query evaluation.}
A critical evaluation detail: the attacker controls query content at test
time (e.g., via a shared trigger deployed to compromised clients).  The
correct evaluation protocol (Chen et al., Algorithm~2) issues triggered
queries $q \oplus \trigger$, not plain queries $q$.  We correct this in our
framework and discuss the resulting performance improvement in
Section~\ref{sec:experiments}.

\paragraph{$\asra$ measurement.}
Beyond modelled $\asra$ (calibrated to paper-reported values,
$\mathcal{N}(0.68, 0.06^2)$), we provide a direct measurement via a local
GPT-2 agent evaluator (\textsc{LocalAgentEvaluator}): retrieved memory
entries are embedded in an agent prompt, GPT-2 generates a response, and
adversarial action execution is detected via keyword matching combined with a
perplexity gate ($\mathrm{ppl} < 200$) that rejects incoherent outputs.
GPT-2 acts as a tractable lower bound on true $\asra$; production LLMs
(GPT-4, Llama-3) are expected to achieve values at or above the modelled
upper bound (Section~\ref{sec:phase17_results}).

\IfFileExists{figures/fig_p17_dpr_convergence.pdf}{%
\begin{figure}[ht]
  \centering
  \includegraphics[width=0.95\linewidth]{fig_p17_dpr_convergence.pdf}
  \caption{DPR HotFlip trigger optimization.
           Left: cosine similarity convergence over iterations (red) compared
           to centroid-passage baseline (orange dashed) and random trigger
           (grey dotted).  Right: summary of trigger quality across methods.
           DPR HotFlip achieves $+0.30$ over random and $+0.07$ over the
           centroid-only approximation.}
  \label{fig:dpr_convergence}
\end{figure}%
}{}

\subsection{MINJA}

\minja{} \citep{dong2025minja} uses a \emph{progressive-shortening indication
protocol}: the attacker sends $T$ crafted queries per target entry, each
progressively shorter to remain below detection thresholds.  The
injection success rate is:
\begin{equation}
  \isr = \frac{|\{\text{entries successfully stored}\}|}{n_\text{poison}},
  \qquad
  \Pr[\text{turn } t \text{ succeeds}] = p_0 \cdot \max(0.5,\, 1 - \lambda t),
\end{equation}
where $p_0 = 0.98$ (full indication probability) and $\lambda = 0.10$
(shortening rate).  With $T=3$ turns, ISR $\approx 0.982$, reproducing the
paper result.

Adversarial passages use \emph{bridging steps} that connect the victim query
topic to the malicious goal via a procedurally correct reasoning chain.
Modelled $\asra \sim \mathcal{N}(0.76, 0.05^2)$ (highest of the three
attacks due to convincing bridging framing).

\subsection{InjecMEM}

\injecmem{} \citep{injecmem2026} employs a retriever-agnostic anchor: a
passage covering broad agent task vocabulary (meetings, tasks, calendar,
security, infrastructure, etc.) such that many query types partially match it.
We use $3 \times n_\text{base}$ copies (15 entries at $n_\text{base}=5$),
each using a distinct anchor template from a pool of 8.  Modelled
$\asra \sim \mathcal{N}(0.57, 0.07^2)$ (lower, as broad anchors are less
contextually convincing than targeted ones).

\subsection{Attack Properties Summary}

Table~\ref{tab:attack_properties} summarises the three threat models across
five dimensions relevant to real-world deployability.

\begin{table}[h]
  \centering
  \small
  \caption{Attack threat-model comparison across five deployment-relevant dimensions.}
  \label{tab:attack_properties}
  \begin{tabular}{lccc}
    \toprule
    Property & \agentpoison{} & \minja{} & \injecmem{} \\
    \midrule
    Memory write access required & \checkmark & --- & --- \\
    Query-only (no write access)  & ---        & \checkmark & \checkmark \\
    Single interaction            & ---        & ---        & \checkmark \\
    Trigger token required        & \checkmark & ---        & --- \\
    Query-specific optimisation   & \checkmark & \checkmark & --- \\
    Retriever-agnostic passage    & ---        & ---        & \checkmark \\
    Persistence (multiple copies) & 1$\times$  & 2$\times$  & 3$\times$ \\
    Modelled $\asra$              & $0.68$ & $0.76$ & $0.57$ \\
    \bottomrule
  \end{tabular}
\end{table}

\subsection{Evaluation Setup}

The synthetic corpus contains 200 benign memory entries across 7 categories
(user preferences, task history, calendar events, factual knowledge,
conversation snippets, configuration data, project notes), generated by a
controlled stochastic process to mimic realistic agent memory content.
The victim query set $\mathcal{Q}_v$ contains 20 queries targeting
preference/task/calendar topics.  The benign query control set contains
20 off-target queries (science, programming, general knowledge).

\paragraph{Encoder choice.}
All retrieval, SAD calibration, and trigger-optimization evaluations use
\texttt{sentence-transformers/all-MiniLM-L6-v2} (384-dim) as the production
retrieval encoder, matching the encoder that a real deployment agent would
use.  The DPR HotFlip trigger optimizer additionally uses the
\texttt{facebook/dpr-ctx\_encoder-single-nq-base} (768-dim BERT-based)
encoder \emph{only} during offline trigger construction (Section~\ref{sec:attacks}),
as the AgentPoison paper specifies DPR for gradient-guided optimization.
This two-encoder setup mirrors the realistic attack scenario in which the
adversary has knowledge of the DPR architecture (used for optimization)
while the deployed system uses a lighter retrieval model (all-MiniLM-L6-v2).
Appendix~\ref{app:ablation} reports ablation results on encoder sensitivity.
